\documentclass[11pt]{article}
\usepackage[a4paper,margin=27mm]{geometry}
\usepackage[T1]{fontenc}
\usepackage{lmodern,microtype}
\usepackage{amsmath,amssymb,amsthm,mathtools,mathrsfs}
\usepackage{booktabs,longtable,array,enumitem,xcolor,xurl}
\usepackage[colorlinks=true,linkcolor=blue!40!black,citecolor=blue!40!black,urlcolor=blue!40!black]{hyperref}
\usepackage{fancyhdr}
\pagestyle{fancy}\fancyhf{}
\fancyhead[L]{\small Conductor descent and distribution ranks}
\fancyhead[R]{\small ProveIt research note}
\fancyfoot[C]{\thepage}
\setlength{\headheight}{14pt}
\setlength{\parskip}{4pt}
\setlength{\parindent}{0pt}
\emergencystretch=3em
\numberwithin{equation}{section}
\newtheorem{theorem}{Theorem}[section]
\newtheorem{proposition}[theorem]{Proposition}
\newtheorem{lemma}[theorem]{Lemma}
\newtheorem{corollary}[theorem]{Corollary}
\newtheorem{conjecture}[theorem]{Conjecture}
\theoremstyle{definition}\newtheorem{definition}[theorem]{Definition}
\newtheorem{example}[theorem]{Example}
\theoremstyle{remark}\newtheorem{remark}[theorem]{Remark}
\DeclareMathOperator{\Li}{Li}
\DeclareMathOperator{\rank}{rank}
\DeclareMathOperator{\ord}{ord}
\DeclareMathOperator{\rad}{rad}
\DeclareMathOperator{\Rea}{Re}
\DeclareMathOperator{\Ima}{Im}
\newcommand{\Q}{\mathbb Q}\newcommand{\C}{\mathbb C}
\newcommand{\Z}{\mathbb Z}\newcommand{\U}{\mathcal U}
\newcommand{\ii}{\mathrm i}
\newcommand{\eps}{\varepsilon}
\newcommand{\ch}{\chi^*}
\newcommand{\bch}{\overline{\chi^*}}
\newcommand{\ds}{\displaystyle}
\newcommand{\Liord}[2]{\mathcal L^{[#1]}(#2)}
\newcommand{\code}[1]{\texttt{\small #1}}
\title{\textbf{Conductor Descent and Flat Distribution Ranks}\\[5pt]
\large Stieltjes jets, cyclotomic polylogarithms,\\and prime-support vanishing}
\author{Research contribution prepared for Vladimir Reshetnikov's ProveIt programme\\
\small Mathematical development and computational assistance: OpenAI ChatGPT}
\date{9 October 2026}
\begin{document}
\maketitle
\thispagestyle{empty}
\begin{abstract}
The Polylogarithms manuscript reports a finite distribution-rank pattern for
parameter derivatives of generalized Stieltjes constants, but does not give a
uniform proof. We define the complete finite-level distribution presentation
and prove a stronger statement: over a characteristic-zero splitting field,
its multiparameter deformation is a free module of rank $\varphi(q)$, with one
explicit generator for every primitive character whose conductor divides $q$.
Consequently the endpoint-fixed rank is $q-\varphi(q)$ for every denominator,
every specialization of the prime weights, and every spectral order. The
result persists for arbitrary jets, with no resonance-induced rank jump.

The same conductor coordinates give finite M\"obius formulas for every
Stieltjes parameter-derivative layer and for character-weighted cyclotomic
polylogarithm traces. In particular, a principal primitive-root trace has an
exact zero of order $\omega(q)-1$ at polylogarithm order one. This produces
explicit logarithmic identities for order derivatives, including a second
order identity at level $30$ and a twisted identity at level $260$.
We also give a small exact certificate showing that the manuscript's
$S_4$ candidate is not a consequence of a specified weight-five
single-product double-shuffle presentation. The candidate itself is not
proved or disproved here. All dimension statements concern presented formal
quotients, not numerical independence of periods.
\end{abstract}
\vspace{4pt}
\textbf{Status.} The theorems below have ordinary mathematical proofs.
Finite matrix and polynomial statements have replayable exact-arithmetic
checks. Floating-point tests are supplementary, not interval certificates.
This note has not been independently refereed or formalized in a proof assistant.
The rank phenomenon is connected to classical universal-distribution theory;
no claim of global priority is made for it or for identities that follow from
standard character and Euler-factor formulas.
\newpage
\setcounter{tocdepth}{1}
\tableofcontents
\newpage

\section{The manuscript target and the scope of the contribution}

The source is the consolidated manuscript \emph{Polylogarithms and their
Arithmetic Bridges} in the ProveIt repository \cite{repo}. Its Chapter~8,
Observation \code{tower:thm:rank}, records an exact finite experiment for
$3\le q\le30$ and parameter-derivative orders $k=1,2,3$:
\begin{equation}\label{eq:target}
(q-1)-\rank=\varphi(q)-1.
\end{equation}
The source explicitly states that a uniform proof is absent and that the
historical scripts were not imported. Chapter~10 lists the divisor-row
structure as a remaining proof problem.

The present work addresses that gap with a fully specified matrix. We do not
identify an unseen historical matrix with our matrix by assumption. Rather,
we prove the rank statement for the natural complete distribution system,
provide its generator and tests, and supply replacement manuscript text that
states the presentation explicitly. Reconciliation with the historical
scripts, should they be recovered, is a separate editorial task.

The decisive observation is that primitive residues at denominator $q$ are
not the most robust coordinates. The right coordinates are character sums
at their \emph{minimal conductors}. In those coordinates, every additional
prime produces a polynomial Euler factor. No division by that factor is
necessary. Thus zeros of Euler factors can destroy a particular coordinate
chart without changing the rank of the distribution module.

Our main outputs are as follows. Theorem~\ref{thm:free} proves a multiparameter
free-module normal form. Corollary~\ref{cor:rank} and
Theorem~\ref{thm:jets} give the uniform rank and jet counts.
Theorem~\ref{thm:stieltjes} gives explicit all-order Stieltjes descent.
Theorems~\ref{thm:polylogdescent} and \ref{thm:principalzero} give the
cyclotomic trace identities and their exact vanishing orders.
Proposition~\ref{prop:s4obstruction} records a distinct finite-presentation
obstruction relevant to the unproved $S_4$ candidate.

\subsection{Relation to existing mathematics}

Kubert proved freeness and primitive-cardinality rank for the universal
ordinary distribution, including its rational representation structure
\cite{kubert}. Ouyang developed a broader universal norm-distribution
framework \cite{ouyang}. Our polynomial deformation belongs to this
mathematical setting. The purpose here is an elementary, explicit
conductor-coordinate derivation tailored to the manuscript's rows, together
with jet formulas and machine-replayable certificates. The general idea of
using characters to diagonalize distributions is not new.

The underlying Hurwitz multiplication law, Laurent expansion and
root-of-unity Fourier formulas are standard \cite{dlmfhurwitz,dlmfpolylog}.
The parameter-derivative ladder is already present in the manuscript and
in Coffey's all-orders work \cite{coffey}; it is not claimed as a new result
of this note. The polylogarithm identities below are derived explicitly,
with special attention to imprimitive characters and the removable
singularity at order one.

\subsection{Two kinds of derivatives}

We use $\zeta^{(n)}(s,a)=\partial_s^n\zeta(s,a)$ for spectral derivatives.
The notation $\gamma_n^{(k)}(a)=\partial_a^k\gamma_n(a)$ denotes a parameter
derivative. To avoid confusion, put
\begin{equation}\label{eq:orderjet}
\Liord{j}{z}:=\left.\frac{\partial^j}{\partial s^j}\Li_s(z)\right|_{s=1}.
\end{equation}
Thus $\Liord{1}{z}$ is an order derivative, not $z\,d\Li_1(z)/dz$ and not
$\Li_0(z)$. This distinction is essential in the new explicit examples.

\section{A precise finite-level distribution presentation}

Fix $q\ge2$. Let $\U_d=(\Z/d\Z)^\times$ for $d>1$, represented by
$1\le a\le d$ with $(a,d)=1$, and put $\U_1=\{1\}$. The fraction $1/1$
represents the endpoint, or residue zero modulo one.

Let $K$ be a characteristic-zero field containing the values of all
Dirichlet characters modulo $q$. Introduce an independent indeterminate
$t_p$ for each prime $p\mid q$ and set
\[
R_q=K[t_p:p\mid q].
\]
Start with the free $R_q$-module on symbols $E(a/d)$, where $d\mid q$ and
$a\in\U_d$. There are exactly $q$ symbols because
$\sum_{d\mid q}\varphi(d)=q$. Fractions are always reduced before symbols
are compared.

For $pd\mid q$, $p$ prime and $a\in\U_d$, impose
\begin{equation}\label{eq:prime-row}
\sum_{j=0}^{p-1}E\!\left(\frac{a+jd}{pd}\right)-t_p E(a/d)=0.
\end{equation}
Call the quotient $\mathscr D_q$. These are formal relations. Substituting
$t_p=p^s$ realizes them by Hurwitz zeta values through
\[
\sum_{j=0}^{p-1}\zeta\!\left(s,\frac{x+j}{p}\right)=p^s\zeta(s,x),
\]
initially for $\Rea s>1$ and then meromorphically.

\begin{definition}[The matrix used in the rank theorem]\label{def:matrix}
Index the $q$ columns by $E(a/q)$, $1\le a\le q$, with fractions reduced
when necessary. For every prime $p\mid q$ and $1\le b\le q/p$, the row is
\begin{equation}\label{eq:matrixrow}
\sum_{j=0}^{p-1}e_{b+jq/p}-t_p e_{pb}.
\end{equation}
The resulting matrix is $M_q(\mathbf t)$. The endpoint-fixed matrix
$\widehat M_q(\mathbf t)$ is obtained by deleting column $q$, not by
deleting rows that involve the endpoint.
\end{definition}

\begin{lemma}[Prime rows suffice]\label{lem:prime}
If $t_d=\prod_{p\mid d}t_p^{v_p(d)}$, every composite distribution row
\[
\sum_{j=0}^{d-1}E((x+j)/d)-t_dE(x)
\]
whose arguments have denominator dividing $q$ is generated by the prime
rows. Conversely, the complete all-divisor presentation includes all prime
rows.
\end{lemma}
\begin{proof}
For $d=uv$, first distribute each of the $u$ preimages through a $v$-fold
cover, then sum. The $uv$ preimages occur once each, while the scalar
coefficient becomes $t_ut_v=t_{uv}$. Subtracting the successive relations
expresses the composite row as a row combination. Induct on the number of
prime factors counted with multiplicity. All intermediate denominators
divide the final denominator, so they remain within the finite level.
\end{proof}

The row with parent argument $1$ is indispensable. For a prime level
$q=\ell$, it is the only row and reads
\begin{equation}\label{eq:primelevel}
\sum_{a=1}^{\ell-1}E(a/\ell)+(1-t_\ell)E(1)=0.
\end{equation}
After fixing $E(1)$, the rank is one. Omitting this row would change the
residual dimension by one and would not model the stated experiment.

\section{A conductor-adapted free basis}

Every character modulo $q$ is induced by a unique primitive character
$\ch$ of conductor $f\mid q$. For $f\mid d\mid q$, write
\begin{equation}\label{eq:Z}
Z_{d,\ch}=\sum_{a\in\U_d}\ch(a)E(a/d).
\end{equation}
At $f=1$ the primitive character is the identically-one character, and
$Z_{1,1}=E(1)$. At other $f$, the primitive character is extended by zero
on integers not coprime to $f$.

Define the polynomial
\begin{equation}\label{eq:Apoly}
A_{d,f,\ch}(\mathbf t)=
\prod_{p\mid f}t_p^{v_p(d)-v_p(f)}
\prod_{\substack{p\mid d\\p\nmid f}}
 t_p^{v_p(d)-1}\bigl(t_p-\ch(p)\bigr).
\end{equation}
An empty product is one.

\begin{theorem}[Polynomial conductor normal form]\label{thm:free}
The module $\mathscr D_q$ is free over $R_q$ of rank $\varphi(q)$.
A basis is
\begin{equation}\label{eq:basis}
\left\{Z_{f,\ch}: f\mid q,\ \ch\text{ primitive of conductor }f\right\}.
\end{equation}
For every $f\mid d\mid q$,
\begin{equation}\label{eq:normal}
Z_{d,\ch}=A_{d,f,\ch}(\mathbf t)Z_{f,\ch}.
\end{equation}
In particular, this normal form remains valid after arbitrary specialization
of the prime weights, including zero weights and roots of unity.
\end{theorem}

\begin{proof}
The group $G=\U_q$ acts by multiplying numerators. Since $K$ contains all
character values and $|G|$ is invertible, the usual character idempotents
split the free module into character components. A denominator-$d$
component contains the character induced by $\ch$ exactly when $f\mid d$,
and then its corresponding eigenspace has rank one, generated by
$Z_{d,\ch}$. Equivalently, finite character orthogonality gives
\begin{equation}\label{eq:inversefourier}
E(a/d)=\frac1{\varphi(d)}
\sum_{\chi\bmod d}\overline{\chi(a)}Z_{d,\chi^*}.
\end{equation}
Here each $\chi$ is replaced on units by its inducing primitive character.

We compute the prime relation in one character component. Sum
\eqref{eq:prime-row} over $a\in\U_d$ with weight $\ch(a)$.
If $p\mid d$, every numerator $a+jd$ is coprime to $pd$, and every unit
modulo $pd$ appears once. Therefore
\begin{equation}\label{eq:edge1}
Z_{pd,\ch}=t_pZ_{d,\ch}\qquad(p\mid d).
\end{equation}
If $p\nmid d$, exactly one of the $p$ numerators is divisible by $p$.
After reducing that fraction, its contribution is
\[
\sum_{b\in\U_d}\ch(pb)E(b/d)=\ch(p)Z_{d,\ch}.
\]
The other numerators are precisely the units modulo $pd$. Hence
\begin{equation}\label{eq:edge2}
Z_{pd,\ch}=(t_p-\ch(p))Z_{d,\ch}\qquad(p\nmid d).
\end{equation}
Every original row is a linear combination of these character-summed rows,
by Fourier inversion on $\U_d$. No character with $f\nmid d$ occurs in
that row module, so no additional relation is missed.

Now fix $\ch$. Its available denominators form the divisor interval
$f\mid d\mid q$. Relations \eqref{eq:edge1}--\eqref{eq:edge2} express every
$Z_{d,\ch}$ as a multiple of $Z_{f,\ch}$ by moving upward from $f$.
The multiplier for a new prime is $t_p-\ch(p)$; every further copy of the
same prime contributes $t_p$. Multipliers from distinct primes commute,
so the result is independent of the chosen path and is exactly
\eqref{eq:Apoly}.

To prove freeness rather than merely spanning, map this character component
to $R_q$ by $Z_{d,\ch}\mapsto A_{d,f,\ch}$. Every edge relation maps to
zero, and the proposed generator maps to one. Thus the component is
isomorphic to $R_q$, with no additional relation on its generator.
Taking the direct sum over characters of $G$ proves freeness of rank
$|G|=\varphi(q)$ and the claimed basis. The construction is polynomial,
so specialization requires no nonvanishing assumption.
\end{proof}

\begin{remark}[What the conductor adds to a character transform]
A character transform only at level $q$ yields $Z_{q,\ch}$, which can
vanish formally when $A_{q,f,\ch}=0$. The generator of the free component
is instead $Z_{f,\ch}$. Forgetting this distinction creates an apparent
rank jump at a singular Euler factor. The normal form proves that the
jump belongs to the chosen coordinates, not to the distribution quotient.
\end{remark}

\begin{corollary}[Uniform distribution rank]\label{cor:rank}
For every specialization $t_p\in\C$,
\begin{equation}\label{eq:allrank}
\rank M_q(\mathbf t)=q-\varphi(q),\qquad
\rank\widehat M_q(\mathbf t)=q-\varphi(q).
\end{equation}
Thus the endpoint-fixed quotient has dimension $\varphi(q)-1$.
For $t_p=p^{k+1}$, the ranks are ranks over $\Q$ and hold for every
integer $k\ge1$ and every $q\ge2$.
\end{corollary}
\begin{proof}
The specialized free quotient has dimension $\varphi(q)$, proving the first
rank formula. Its endpoint $E(1)=Z_{1,1}$ is one of the conductor basis
vectors. Thus quotienting further by $E(1)$ removes exactly one dimension.
Equivalently, projection deleting the endpoint column is injective on the
row space: a row-space element supported only on the endpoint would give
a nonzero relation on this basis vector. This proves the second formula.
For integer $k$, all entries are rational, and extending the coefficient
field to $\C$ does not change matrix rank.
\end{proof}

\subsection{A fully explicit polynomial example at level twelve}

Let $\chi_3$ be the nontrivial character modulo $3$, $\chi_4$ the
nontrivial character modulo $4$, and $\chi_{12}=\chi_3\chi_4$.
The four conductor generators are
\begin{align}
Z_1&=E(1),&
Z_3&=E(1/3)-E(2/3),\\
Z_4&=E(1/4)-E(3/4),&
Z_{12}&=E(1/12)-E(5/12)-E(7/12)+E(11/12).
\end{align}
For example,
\begin{align}\label{eq:q12example}
E(1/12)=\frac14\bigl[&t_2(t_2-1)(t_3-1)Z_1
+t_2(t_2+1)Z_3\notag\\
&+(t_3+1)Z_4+Z_{12}\bigr],\\
E(1/6)&=\frac12\bigl[(t_2-1)(t_3-1)Z_1+(t_2+1)Z_3\bigr],\\
E(1/4)&=\frac12\bigl[t_2(t_2-1)Z_1+Z_4\bigr].
\end{align}
The remaining rows follow by character signs and denominator reduction.
The artifact \code{q12\_polynomial\_normal\_form.json} contains the entire
$12\times4$ matrix $F$ and the $4\times12$ extraction matrix $B$.
The exact checks are
\begin{equation}\label{eq:certificateq12}
M_{12}F=0,\qquad BF=I_4.
\end{equation}
They hold as polynomial identities in independent $t_2,t_3$, not merely
at numerical weights.

\section{Jets and the uniform Stieltjes rank theorem}

Let $s_0\in\C$ and $N\ge0$. For a family of symbols $X_a(s)$ satisfying
Hurwitz-type distribution, introduce divided Taylor coefficients
$X_{a,n}=[\eps^n]X_a(s_0+\eps)$. The prime relations are
\begin{equation}\label{eq:jetrows}
\sum_{j=0}^{p-1}X_{b+jq/p,n}
-p^{s_0}\sum_{r=0}^{n}\frac{(\log p)^r}{r!}X_{pb,n-r}=0,
\qquad 0\le n\le N.
\end{equation}
All powers use the real logarithm of the positive integer $p$.

\begin{theorem}[No rank loss under arbitrary jets]\label{thm:jets}
The complete system \eqref{eq:jetrows} on $(N+1)q$ coordinates has rank
\begin{equation}\label{eq:jetrank}
(N+1)\bigl(q-\varphi(q)\bigr).
\end{equation}
Its solution space has dimension $(N+1)\varphi(q)$. Fixing all endpoint
jets leaves $(N+1)(\varphi(q)-1)$ free coordinates and does not change
the row rank. These statements hold for every complex $s_0$.
\end{theorem}
\begin{proof}
Work with truncated germs modulo $\eps^{N+1}$ and substitute
$t_p=p^{s_0}\exp(\eps\log p)$ into Theorem~\ref{thm:free}.
Equivalently, choose arbitrary truncated germs for the conductor-character
coordinates and define all other character germs by
\eqref{eq:normal}; then use Fourier inversion. This produces every solution
of \eqref{eq:jetrows}, because the character-summed equations still have
unit coefficient on the higher-denominator coordinate. Extracting the
minimal-conductor sums recovers the chosen germs, proving both completeness
and uniqueness. There are $\varphi(q)$ freely chosen germs, each of length
$N+1$. The endpoint is precisely the principal-conductor germ and can be
fixed independently. Subtracting the solution-space dimension from the
number of columns gives the rank.
\end{proof}

The proof also applies when the quantities replacing $\log p$ are arbitrary
complex numbers, or independent formal variables. In particular, no
algebraic independence of logarithms is being used.

The generalized Stieltjes constants are defined by
\begin{equation}\label{eq:Laurent}
\zeta(1+\eps,a)=\frac1\eps+
\sum_{n\ge0}\frac{(-1)^n}{n!}\gamma_n(a)\eps^n.
\end{equation}
For $k\ge1$, differentiating in $a$ removes the pole and gives the
known generating ladder
\begin{equation}\label{eq:master}
\sum_{n\ge0}\frac{(-1)^n}{n!}\gamma_n^{(k)}(a)\eps^n
=(-1)^k(1+\eps)_k\zeta(k+1+\eps,a).
\end{equation}
For completeness, $\partial_a^k\zeta(s,a)=(-1)^k(s)_k\zeta(s+k,a)$
follows by termwise differentiation for $\Rea s>1$ and then continuation.
Substitution into \eqref{eq:Laurent} proves \eqref{eq:master}.

\begin{corollary}[The manuscript's rank pattern, uniformly]\label{cor:stielrank}
Fix $q\ge2$, $k\ge1$ and $n\ge0$. Treat endpoint values and lower
Stieltjes indices as known. The complete distribution system on
\[
\gamma_n^{(k)}(a/q),\qquad 1\le a<q,
\]
has coefficient rank $q-\varphi(q)$ and formal residual dimension
$\varphi(q)-1$. If indices $0\le n\le N$ are retained together, the
residual dimension is $(N+1)(\varphi(q)-1)$.
\end{corollary}
\begin{proof}
Multiplication by $(1+\eps)_k$ is invertible on truncated germs since its
constant coefficient is $k!\ne0$. It is common to every argument and
therefore preserves distribution. The top-index coefficient matrix is
$\widehat M_q$ with $t_p=p^{k+1}$. Lower-index logarithmic terms belong to
the inhomogeneous right-hand side. Apply Corollary~\ref{cor:rank} and
Theorem~\ref{thm:jets}.
\end{proof}

This proves a uniform statement about the explicit law presentation. It
does not prove that the evaluated constants occupy that many independent
numerical directions. Analytic functional equations or unrecorded period
relations may reduce the evaluated span further.

\section{Finite conductor descent for all Stieltjes layers}

For a primitive character $\ch$ of conductor $f\mid d$, define
\begin{equation}
G_{d,\ch}^{(n,k)}=
\sum_{a\in\U_d}\ch(a)\gamma_n^{(k)}(a/d).
\end{equation}
Set
\begin{equation}\label{eq:Rdf}
R(d,f)=\prod_{\substack{p\mid d\\p\nmid f}}p,
\qquad m_r=\frac{d}{fr}\quad(r\mid R(d,f)).
\end{equation}
Each $m_r$ is a positive integer. Upon substituting $t_p=p^s$,
\eqref{eq:Apoly} becomes the entire exponential polynomial
\begin{align}\label{eq:Aanalytic}
A_{d,f,\ch}(s)
&=\left(\frac df\right)^s
\prod_{\substack{p\mid d\\p\nmid f}}(1-\ch(p)p^{-s})\notag\\
&=\sum_{r\mid R(d,f)}\mu(r)\ch(r)m_r^s.
\end{align}
In particular,
\begin{equation}\label{eq:Aderivative}
A_{d,f,\ch}^{(j)}(s)
=\sum_{r\mid R(d,f)}\mu(r)\ch(r)m_r^s\log^j m_r.
\end{equation}

\begin{theorem}[All-index, all-parameter-order descent]\label{thm:stieltjes}
For $n\ge0$, $k\ge1$ and $f\mid d$,
\begin{align}\label{eq:stieldescent}
G_{d,\ch}^{(n,k)}
={}&\sum_{r\mid R(d,f)}\mu(r)\ch(r)m_r^{k+1}\notag\\
&\quad\cdot\sum_{j=0}^{n}(-1)^j\binom nj\log^j m_r\,
G_{f,\ch}^{(n-j,k)}.
\end{align}
No Euler-factor denominator occurs. Equivalently,
\begin{equation}\label{eq:stielA}
G_{d,\ch}^{(n,k)}=
\sum_{j=0}^{n}(-1)^j\binom nj A_{d,f,\ch}^{(j)}(k+1)
G_{f,\ch}^{(n-j,k)}.
\end{equation}
\end{theorem}
\begin{proof}
For $\Rea s>1$, sorting a Dirichlet series by residue classes gives
\[
\sum_{a\in\U_d}\ch(a)\zeta(s,a/d)=d^sL(s,\chi_d),
\]
where $\chi_d$ is the induced character modulo $d$. Removing the extra
Euler factors from its primitive $L$-series yields
\[
L(s,\chi_d)=L(s,\ch)
\prod_{\substack{p\mid d\\p\nmid f}}(1-\ch(p)p^{-s}).
\]
Thus the character sum at $d$ is $A_{d,f,\ch}(s)$ times the character
sum at $f$. This is also the analytic realization of
Theorem~\ref{thm:free}.

Insert $s=k+1+\eps$, multiply by the common factor
$(-1)^k(1+\eps)_k$, and use \eqref{eq:master}. Comparison of
$(-1)^n\eps^n/n!$ gives \eqref{eq:stielA}; the sign $(-1)^j$ comes
from the Stieltjes Laurent convention. Finally apply
\eqref{eq:Aderivative}.
\end{proof}

\subsection{Concrete second-index identities at level twelve}

The multipliers are
\[
A_{12,3,\chi_3}(s)=4^s+2^s,\qquad
A_{12,4,\chi_4}(s)=3^s+1.
\]
Therefore, at $n=2$, $k=1$,
\begin{align}\label{eq:examplegamma}
G_{12,\chi_3}^{(2,1)}
&=20G_{3,\chi_3}^{(2,1)}
-72\log2\,G_{3,\chi_3}^{(1,1)}
+68\log^22\,G_{3,\chi_3}^{(0,1)},\\
G_{12,\chi_4}^{(2,1)}
&=10G_{4,\chi_4}^{(2,1)}
-18\log3\,G_{4,\chi_4}^{(1,1)}
+9\log^23\,G_{4,\chi_4}^{(0,1)}.
\end{align}
For instance,
\[
G_{12,\chi_3}^{(2,1)}=
\gamma_2'(1/12)-\gamma_2'(5/12)+\gamma_2'(7/12)-\gamma_2'(11/12),
\]
and $G_{3,\chi_3}^{(2,1)}=\gamma_2'(1/3)-\gamma_2'(2/3)$.
These are exact identities between explicitly specified signed sums; they
require no integer-relation search. The two identities were also checked
numerically using independent finite Hurwitz character sums.

\subsection{The undifferentiated layer and its pole correction}

For $k=0$, a principal character sum still contains the Hurwitz pole.
The same formula without a correction would be wrong.

\begin{proposition}[The $k=0$ correction]\label{prop:kzero}
For nonprincipal primitive $\ch$ the formula
\eqref{eq:stielA} remains valid with $k=0$. For $f=1$, put
$A_d(s)=A_{d,1,1}(s)$. Then
\begin{equation}\label{eq:kzero}
G_{d,1}^{(n,0)}=
\sum_{j=0}^{n}(-1)^j\binom nj A_d^{(j)}(1)\gamma_{n-j}(1)
+\frac{(-1)^n}{n+1}A_d^{(n+1)}(1).
\end{equation}
\end{proposition}
\begin{proof}
For a nonprincipal character, the sum of its values over units is zero,
so the $1/\eps$ residues cancel. Coefficient comparison is unchanged.
For the principal conductor, multiply
\[
A_d(1+\eps)\left(\frac1\eps+
\sum_{n\ge0}\frac{(-1)^n\gamma_n(1)}{n!}\eps^n\right).
\]
The pole term contributes $A_d^{(n+1)}(1)/(n+1)!$ to the coefficient of
$\eps^n$. Multiplying by $(-1)^n n!$ gives the correction in
\eqref{eq:kzero}. The residue itself is
$A_d(1)=d\prod_{p\mid d}(1-p^{-1})=\varphi(d)$, as required.
\end{proof}

For the non-character multiplication formula this specializes to
\begin{equation}\label{eq:classicalstieldist}
\sum_{j=0}^{d-1}\gamma_n\!\left(\frac{a+j}{d}\right)
=d\sum_{r=0}^{n}\binom nr(-\log d)^r\gamma_{n-r}(a)
+\frac{(-1)^n d\log^{n+1}d}{n+1}.
\end{equation}
The inhomogeneous term is the residue contribution, not an optional
normalization.

\section{Cyclotomic polylogarithm conductor descent}

Put $\xi_d=e^{2\pi\ii/d}$. For $f\mid d$ and a primitive character
$\ch$ of conductor $f$, define the primitive-root trace
\begin{equation}\label{eq:Pdef}
P_{d,\ch}(s)=\sum_{a\in\U_d}\bch(a)\Li_s(\xi_d^a).
\end{equation}
At $d=1$, the trace is $\zeta(s)$. At every $d>1$, all roots in the sum
are different from one, and the trace is entire in the order $s$.
One proof of this last assertion is the finite Fourier formula
\begin{equation}\label{eq:fourierzeta}
\Li_s(\xi_d^a)=d^{-s}\sum_{b=1}^{d}\xi_d^{ab}\zeta(s,b/d),
\qquad a\not\equiv0\pmod d.
\end{equation}
The potential pole at $s=1$ cancels because $\sum_b\xi_d^{ab}=0$.

Let
\begin{equation}\label{eq:B}
B_{d,f,\ch}(s)=
\sum_{r\mid R(d,f)}\mu(r)\bch(r)
\left(\frac{d}{fr}\right)^{1-s}
=A_{d,f,\overline{\ch}}(1-s).
\end{equation}
The conjugation and the reflection $s\mapsto1-s$ are both important.

\begin{theorem}[Primitive-root trace descent]\label{thm:polylogdescent}
As a meromorphic identity in $s$,
\begin{equation}\label{eq:Pdescent}
P_{d,\ch}(s)=B_{d,f,\ch}(s)P_{f,\ch}(s).
\end{equation}
For $f>1$ both sides are entire. For $f=1<d$, the right-hand side has a
removable singularity at one and equals the entire left-hand side.
For a primitive character,
\begin{equation}\label{eq:gauss}
P_{f,\ch}(s)=\tau(\bch)L(s,\ch),\qquad
\tau(\bch)=\sum_{a\in\U_f}\bch(a)\xi_f^a.
\end{equation}
\end{theorem}

\begin{proof}
It suffices first to work in $\Rea s>1$, where all Dirichlet series are
absolutely convergent. Write $d=fm$. In the trace defining $P_d$, the
restriction $(a,d)=1$ can be imposed, after requiring $(a,f)=1$, by
inclusion--exclusion over $r\mid R(d,f)$. Put $a=rb$ in the term indexed
by $r$. The inner exponential sum in the coefficient of $n^{-s}$ is
\[
\bch(r)\sum_{b=1}^{d/r}\bch(b)e^{2\pi\ii bn/(d/r)}.
\]
Split $b$ into classes modulo $f$. The sum vanishes unless $(m/r)\mid n$.
When $n=(m/r)\ell$, it equals
\[
\bch(r)\frac mr
\sum_{b\in\U_f}\bch(b)e^{2\pi\ii b\ell/f}.
\]
For a primitive character, the last sum is
$\tau(\bch)\ch(\ell)$, including zero when $(\ell,f)>1$.
Summing over $n$ gives
\[
P_{d,\ch}(s)=\tau(\bch)L(s,\ch)
\sum_{r\mid R(d,f)}\mu(r)\bch(r)(m/r)^{1-s}.
\]
The special case $d=f$ gives \eqref{eq:gauss}, and division is not needed
to obtain \eqref{eq:Pdescent}. Analytic continuation and
\eqref{eq:fourierzeta} finish the proof. The primitive Gauss identity used
here is the standard finite Fourier identity for a primitive character;
it follows by permuting units when $(\ell,f)=1$, with primitivity forcing
vanishing at nonunits.
\end{proof}

\begin{corollary}[All order derivatives away from the principal pole]
If $f>1$, then for every $j\ge0$ and every $s_0\in\C$,
\begin{align}\label{eq:Pjets}
P_{d,\ch}^{(j)}(s_0)
={}&\sum_{r\mid R(d,f)}\mu(r)\bch(r)m_r^{1-s_0}\notag\\
&\quad\cdot\sum_{h=0}^{j}\binom jh(-\log m_r)^h
P_{f,\ch}^{(j-h)}(s_0).
\end{align}
For $f=1$ the same formula holds away from $s_0=1$; at one, its finite
part must be taken before evaluation.
\end{corollary}
\begin{proof}
Differentiate the finite exponential sum in \eqref{eq:Pdescent}. For
$f>1$ the factors are entire, so ordinary Leibniz differentiation applies
at every order.
\end{proof}

\begin{example}[Two nontrivial level-twelve identities]
For all complex $s$,
\begin{align}\label{eq:P12}
P_{12,\chi_4}(s)&=(1+3^{1-s})P_{4,\chi_4}(s),\\
P_{12,\chi_3}(s)&=(4^{1-s}+2^{1-s})P_{3,\chi_3}(s).
\end{align}
The first identity is
\begin{align*}
&\Li_s(\xi_{12})+\Li_s(\xi_{12}^{5})
-\Li_s(\xi_{12}^{7})-\Li_s(\xi_{12}^{11})\\
&\hspace{25mm}=(1+3^{1-s})\bigl(\Li_s(\ii)-\Li_s(-\ii)\bigr).
\end{align*}
These were checked directly at the nonintegral complex order
$s=1.7+0.3\ii$, separately from the proof.
\end{example}

\section{Exact prime-support vanishing at order one}

The principal trace is particularly transparent. Put
\begin{equation}\label{eq:Tdef}
T_q(s)=\sum_{a\in\U_q}\Li_s(\xi_q^a),\qquad q>1.
\end{equation}
Theorem~\ref{thm:polylogdescent} gives
\begin{equation}\label{eq:Tprod}
T_q(s)=q^{1-s}\prod_{p\mid q}(1-p^{s-1})\zeta(s).
\end{equation}
This is also the Dirichlet series of the classical Ramanujan sum
$c_q(n)=\sum_{a\in\U_q}\xi_q^{an}$, obtained directly by
inclusion--exclusion. We use the product to examine the removable
singularity, rather than substituting $s=1$ into the factors separately.

\begin{theorem}[Exact vanishing order]\label{thm:principalzero}
Let $r=\omega(q)$ be the number of distinct prime divisors of $q>1$.
Then the entire function $T_q$ has an exact zero of order $r-1$ at $s=1$.
Here order zero means a nonzero value. Explicitly,
\begin{align}\label{eq:zeros}
\sum_{a\in\U_q}\Liord{j}{\xi_q^a}&=0&& (0\le j<r-1),\\
\sum_{a\in\U_q}\Liord{r-1}{\xi_q^a}
&=(-1)^r(r-1)!\prod_{p\mid q}\log p.\label{eq:leading}
\end{align}
\end{theorem}
\begin{proof}
Set $s=1+\eps$. Every factor
$1-p^\eps=-\eps\log p+O(\eps^2)$ has a simple zero, while
$\zeta(1+\eps)=\eps^{-1}+O(1)$ has a simple pole.
The nonzero factor $q^{-\eps}$ equals one at zero. Therefore
\begin{equation}\label{eq:Tleading}
T_q(1+\eps)=(-1)^r\left(\prod_{p\mid q}\log p\right)
\eps^{r-1}+O(\eps^r).
\end{equation}
The leading coefficient is nonzero because every prime exceeds one.
Taylor coefficient extraction proves both statements.
\end{proof}

\subsection{Explicit identities}

At level $q=30$, $\U_{30}=\{1,7,11,13,17,19,23,29\}$, and
\begin{align}\label{eq:q30}
\sum_{a\in\U_{30}}\Li_1(\xi_{30}^a)&=0,\\
\sum_{a\in\U_{30}}\Liord{1}{\xi_{30}^a}&=0,\\
\sum_{a\in\U_{30}}\Liord{2}{\xi_{30}^a}
&=-2\log2\log3\log5.
\end{align}
At level $210$, the first three traces, of derivative orders $0,1,2$,
vanish, and
\begin{equation}\label{eq:q210}
\sum_{a\in\U_{210}}\Liord{3}{\xi_{210}^a}
=6\log2\log3\log5\log7.
\end{equation}
Prime powers do not change the first nonzero derivative: at $q=p^e$,
$T_q(1)=-\log p$. Their exponents do affect later derivatives.

\subsection{The next two coefficients, including repeated prime powers}

Write
\[
L=\log q,\quad S_1=\sum_{p\mid q}\log p,\quad
S_2=\sum_{p\mid q}(\log p)^2,
\quad C=(-1)^r\prod_{p\mid q}\log p,
\]
and let $\gamma=\gamma_0(1)$, $\gamma_1=\gamma_1(1)$.
Expanding \eqref{eq:Tprod} two further orders gives
\begin{align}\label{eq:Tnext}
T_q^{(r)}(1)
&=r!C\left(\gamma-L+\frac{S_1}{2}\right),\\
T_q^{(r+1)}(1)
&=(r+1)!C\left[-\gamma_1+
\gamma\left(\frac{S_1}{2}-L\right)
+\frac12\left(\frac{S_1}{2}-L\right)^2+\frac{S_2}{24}\right].
\end{align}
To verify the expansion, use
\[
\frac{e^x-1}{x}=1+\frac x2+\frac{x^2}{6}+O(x^3),\qquad
\log\frac{e^x-1}{x}=\frac x2+\frac{x^2}{24}+O(x^3),
\]
and $\eps\zeta(1+\eps)=1+\gamma\eps-\gamma_1\eps^2+O(\eps^3)$.
This gives an explicit all-prime-support family beyond the leading
logarithmic identity. Higher coefficients are obtained by finite series
multiplication with the same sign convention.

\section{Twisted vanishing and a level-260 identity}

Let $\ch$ be primitive of conductor $f>1$, and $f\mid q$.
Among the added primes, define
\[
\mathcal S=\{p\mid q:p\nmid f,\ \ch(p)=1\},\qquad r_\chi=|\mathcal S|.
\]
The remaining added primes form $\mathcal N$.

\begin{theorem}[Twisted leading jet]\label{thm:twisted}
The trace $P_{q,\ch}$ vanishes to order at least $r_\chi$ at one, and
\begin{equation}\label{eq:twistedlead}
P_{q,\ch}^{(r_\chi)}(1)=
(-1)^{r_\chi}r_\chi!\,
P_{f,\ch}(1)
\prod_{p\in\mathcal S}\log p
\prod_{p\in\mathcal N}(1-\bch(p)).
\end{equation}
Its order is exactly $r_\chi$ whenever $P_{f,\ch}(1)\ne0$.
More generally its exact order is $r_\chi+\ord_{s=1}P_{f,\ch}(s)$.
\end{theorem}
\begin{proof}
The multiplier is
\[
B_{q,f,\ch}(1+\eps)=
(q/f)^{-\eps}
\prod_{\substack{p\mid q\\p\nmid f}}(1-\bch(p)p^\eps).
\]
The factors indexed by $\mathcal S$ have simple zeros with leading
coefficient $-\log p$. All other factors are nonzero at zero and have
value $1-\bch(p)$. Multiply their leading terms by the Taylor series
of $P_{f,\ch}(1+\eps)$.
\end{proof}

For primitive nonprincipal characters, the classical nonvanishing
$L(1,\ch)\ne0$, together with the nonzero primitive Gauss sum in
\eqref{eq:gauss}, supplies the stated nonvanishing hypothesis
\cite{dlmfdirichlet}. The explicit example below does not need that theorem;
its primitive trace is computed directly from logarithms.

For $\chi_4(a)=(-1)^{(a-1)/2}$ on odd integers,
\[
P_{4,\chi_4}(1)=\Li_1(\ii)-\Li_1(-\ii)=\frac{\ii\pi}{2}.
\]
Take $q=260=4\cdot5\cdot13$. Both added primes have $\chi_4(p)=1$.
Consequently
\begin{align}\label{eq:260}
\sum_{a\in\U_{260}}\chi_4(a)\Li_1(\xi_{260}^a)&=0,\\
\sum_{a\in\U_{260}}\chi_4(a)\Liord{1}{\xi_{260}^a}&=0,\\
\sum_{a\in\U_{260}}\chi_4(a)\Liord{2}{\xi_{260}^a}
&=\ii\pi\log5\log13.
\end{align}
This is a proved identity, not an integer-relation conjecture. More generally,
for distinct primes $p_1,\ldots,p_r\equiv1\pmod4$ and
$q=4p_1\cdots p_r$,
\begin{equation}\label{eq:chi4family}
\sum_{a\in\U_q}\chi_4(a)\Liord{r}{\xi_q^a}
=(-1)^r r!\frac{\ii\pi}{2}\prod_{j=1}^r\log p_j,
\end{equation}
with all lower order traces vanishing. The same leading identity persists
if the added prime factors occur with higher exponents.

\section{Evaluation without differentiating through a removable pole}

The trace identities concern an analytic function at $s=1$, but a numerical
implementation of $\Li_s$ may internally use expressions whose gamma and
zeta factors cancel only after substantial subtraction. Tiny-step
black-box differentiation can then amplify roundoff.

There is an exact alternative. Expand \eqref{eq:fourierzeta} using
\eqref{eq:Laurent} and remove the pole \emph{symbolically} before extracting
coefficients.

\begin{proposition}[Pole-subtracted Fourier jet formula]\label{prop:eval}
For $q\ge2$, $a\not\equiv0\pmod q$ and $n\ge0$,
\begin{equation}\label{eq:Fourierjet}
\Liord{n}{\xi_q^a}
=\frac{(-1)^n}{q}\sum_{b=1}^{q}\xi_q^{ab}
\sum_{j=0}^{n}\binom nj(\log q)^{n-j}\gamma_j(b/q).
\end{equation}
\end{proposition}
\begin{proof}
The pole coefficient is proportional to $\sum_b\xi_q^{ab}=0$.
The remaining product is
\[
q^{-1}e^{-\eps\log q}
\sum_{b=1}^q\xi_q^{ab}
\sum_{j\ge0}\frac{(-1)^j}{j!}\gamma_j(b/q)\eps^j.
\]
Its $n$th Taylor coefficient, multiplied by $n!$, is
\eqref{eq:Fourierjet}.
\end{proof}

This formula is especially efficient for a whole primitive-root trace:
precompute the Stieltjes table once and sum with its finite Fourier weights.
For the principal trace those weights are the integer Ramanujan sums
$c_q(b)$. It also avoids pretending that a floating-point sum of roots is
\emph{exactly} zero before multiplying by a numerical pole.

In the local audit, a direct tiny-step order differentiation of
\code{mpmath.polylog} near order one produced a spurious approximately
$1.4\cdot10^{-9}$ first trace at $q=30$. The final evaluator did not loosen
the acceptance tolerance. It instead used \eqref{eq:Fourierjet}; all ten
final smoke tests at 45 decimal working digits passed the $10^{-38}$
absolute-error threshold. The largest observed error was approximately
$3.59\cdot10^{-43}$. This is an implementation observation for the tested
version and settings, not a theorem about every differentiation routine.

An interval implementation remains desirable. The current numerical
artifacts explicitly label themselves as floating-point smoke tests. The
analytic identities are established by the proofs above, independently of
those residuals.

\section{An exact limitation of the tested laws for the \texorpdfstring{$S_4$}{S4} candidate}

The manuscript defines
\[
S_4=\sum_{n\ge0}\frac{(-1)^nH_n}{(2n+1)^4},\qquad
g_{a,b}=\Ima\Li_{a,b}(\ii,1),
\]
and displays the candidate
\begin{equation}\label{eq:S4candidate}
S_4=\frac{4g_{4,1}-3g_{3,2}-9g_{2,3}}7
+\frac{\pi^5}{224}-\frac{27G\zeta(3)}{224}
-2\beta(4)\log2,
\end{equation}
where $G=\beta(2)$. We do not promote this identity to theorem status.
We do, however, explain one precise obstacle to proving it from a restricted
set of laws.

\begin{lemma}[The colored double represented by $S_p$]\label{lem:Sdouble}
For $p>1$,
\begin{equation}\label{eq:Sdouble}
S_p=\Ima\left(\Li_{p,1}(\ii,1)+\Li_{p,1}(\ii,-1)\right).
\end{equation}
\end{lemma}
\begin{proof}
In the sum of the two doubles the inner color factor is $1+(-1)^m$,
which retains twice the even indices. At outer index $2n+1$, the inner
sum is $2\sum_{j=1}^n(2j)^{-1}=H_n$, while
$\Ima(\ii^{2n+1})=(-1)^n$. Even outer indices contribute no imaginary
part. Absolute convergence for $p>1$ justifies the regrouping.
\end{proof}

At weight five put
\[
x_{a,b;c,d}=\Ima\Li_{a,b}(\ii^c,\ii^d),\qquad a+b=5,
\]
formally imposing conjugation
$x_{a,b;-c,-d}=-x_{a,b;c,d}$. When both colors are real, the coordinate is
zero. The six representative color pairs are
\[
(0,1),(1,0),(1,1),(1,2),(1,3),(2,1).
\]
There are thus $24$ coordinates. In this paragraph all products of single
polylogarithms and single weight-five values are quotiented out.

The tested presentation consists of the imaginary parts of the following
two formal product rows, for $p+q=5$ and $c,d\in\Z/4\Z$:
\begin{align}\label{eq:Srows}
&x_{p,q;c,d}+x_{q,p;d,c}=0,\\
&\sum_{j=0}^{p-1}\binom{q-1+j}{j}x_{q+j,p-j;d,c-d}
+\sum_{j=0}^{q-1}\binom{p-1+j}{j}x_{p+j,q-j;c,d-c}=0.
\end{align}
They are the depth-one-product stuffle and shuffle coefficient rows.
The formal presentation includes the endpoint-color rows as well; restricting
to convergent products alone can only reduce its row space. Nonzero rows,
including duplicates, form a $96\times24$ integer matrix $M$.
The ordinary level-two distribution sums of these color families have
zero imaginary row after conjugation, so adding those rows does not alter
this matrix rank.

By Lemma~\ref{lem:Sdouble}, the depth-two left side of seven times
\eqref{eq:S4candidate} is
\begin{equation}\label{eq:S4target}
\ell=7x_{4,1;1,2}+3x_{4,1;1,0}
+3x_{3,2;1,0}+9x_{2,3;1,0}.
\end{equation}
Its conjectured single-value right side is
$\pi^5/32-27G\zeta(3)/32-14\beta(4)\log2$.

\begin{proposition}[A finite-presentation obstruction]\label{prop:s4obstruction}
The row $\ell$ is not in the row space of the presentation
\eqref{eq:Srows}. In fact,
\[
\rank M=20,\qquad \rank\begin{pmatrix}M\\\ell\end{pmatrix}=21.
\]
Thus \eqref{eq:S4candidate} cannot follow merely by linear combination
of the tested same-weight single-product rows, after the specified
single-value quotient. This is not a disproof of the candidate.
\end{proposition}
\begin{proof}
It suffices for the nonmembership assertion to give a column vector $v$
annihilated by all rows of $M$ but not by $\ell$. Its only nonzero entries
are the following:
\begin{center}
\begin{tabular}{ccccr}
\toprule
$a$&$b$&$c$&$d$&$v_{a,b;c,d}$\\\midrule
1&4&1&0&$-2$\\
2&3&0&1&$1$\\
2&3&1&0&$3$\\
2&3&1&3&$-1$\\
3&2&0&1&$-3$\\
3&2&1&0&$-1$\\
3&2&1&3&$-1$\\
4&1&0&1&$2$\\\bottomrule
\end{tabular}
\end{center}
Substitution in the finite binomial sums \eqref{eq:Srows}, with the stated
conjugation rule, gives $Mv=0$. On the other hand,
$\ell v=3(-1)+9(3)=24$. Hence $\ell$ is not a row combination.
The two displayed ranks are separately verified by exact integer/rational
row reduction in \code{verify\_s4\_obstruction.py}. The full matrix, row
labels and witness are retained in the JSON certificate; no numerical
values of polylogarithms are used in this check.
\end{proof}

\subsection{What remains plausible and what was numerically checked}

The conjecture was independently tested using the one-dimensional integrals
\begin{align}\label{eq:S4int}
S_4&=-\frac16\int_0^1\frac{(-\log t)^3\log(1+t^2)}{1+t^2}\,dt,\\
g_{a,b}&=\frac1{(a-1)!}\Ima\int_0^1
(-\log t)^{a-1}\frac{\ii\Li_b(\ii t)}{1-\ii t}\,dt.
\end{align}
The first follows from the harmonic-number generating function; the second
from the derivative of $\Li_{1,b}(z,1)$ and repeated logarithmic
integration. An initial 50-digit check gave residual magnitude about
$2.32\cdot10^{-51}$; the packaged 45-digit rerun also passed.
This is evidence, not a proof.

The exact obstruction says where the next proof attempt must look: an
additional functional equation, a higher-depth relation whose elimination
creates a new double relation, or an argument transformation outside the
single-product presentation. Enlarging the coefficient height in the same
law span cannot solve the problem. It would be incorrect either to reject
\eqref{eq:S4candidate} because of this obstruction or to claim it is proved
because its floating-point residual is small.

\section{Corrections and integration recommendations}

\subsection{Replace the finite rank observation by a precisely scoped theorem}

The most substantial editorial change is to replace
\code{tower:thm:rank} with Corollary~\ref{cor:stielrank}, including
Definition~\ref{def:matrix}. State that the theorem concerns the complete
prime/distribution presentation with the endpoint fixed. Do not replace it
by a bare numerical-independence assertion. The standalone replacement
fragment is supplied in the integration directory.

\subsection{Make the weight-one endpoint explicit}

Some prose surrounding the Gaussian and Eisenstein single values says that
their real parts are rational multiples of $\zeta(n)$. The $n=1$ case must
be separated or interpreted by a removable limit. In the principal branch,
\begin{align}\label{eq:Li1}
\Li_1(\ii)&=-\tfrac12\log2+\tfrac{\ii\pi}{4},\\
\Li_1(e^{2\pi\ii/3})&=-\tfrac12\log3+\tfrac{\ii\pi}{6}.
\end{align}
For example, $\tfrac12(3^{1-s}-1)\zeta(s)$ tends to
$-\tfrac12\log3$ at $s=1$; literal substitution gives an undefined
zero times pole. This is a domain clarification, not evidence that the
higher-weight formulas are wrong.

\subsection{Do not invert a vanishing conductor multiplier}

A formula that solves for $Z_{f,\ch}$ by dividing $Z_{q,\ch}$ by
$A_{q,f,\ch}(s)$ needs a nonzero-denominator hypothesis. The polynomial
normal form avoids it. The same warning applies to the polylogarithm
multiplier $B_{q,f,\ch}$, where the zeros at one produce the identities in
Sections~8 and~9. A disappearing primitive-root trace is not a disappearing
formal character direction.

\subsection{Keep the undifferentiated pole term}

Any extension of the manuscript's $k\ge1$ identities to $k=0$ must include
Proposition~\ref{prop:kzero} for the principal character. The extra term is
forced by the Laurent pole. Its omission would already fail at $n=0$.

\subsection{Preserve the status of the remaining \texorpdfstring{$S_4$}{S4} candidate}

The earlier weight-five sporadic relation and weight-six Gaussian parity
reductions are not re-presented here as new work. The displayed $S_4$
formula remains a candidate in this contribution. Add its exact
single-product obstruction as a methodological note, not as a negative
numerical claim. A later proof using a larger law vocabulary would be
fully compatible with the certificate given here.

\section{Reproducibility, verification and implementation boundaries}

The package contains three kinds of evidence. First, this article proves
the general statements without relying on finite experiments. Second,
exact scripts verify finite matrices and a polynomial normal form. Third,
independent high-precision numerical representations test signs,
conjugations and coefficient conventions.

The completed exact run contains 295 specialized rank cases: $2\le q\le60$
at weights $t_p=1,p^2,p^3$, at all-zero prime weights, and at one mixed
integer assignment extended multiplicatively. Each checks both the full
matrix and its endpoint-fixed projection. There are 29 all-divisor checks
for $2\le q\le30$, and 18 symbolic jet cases at
$q\in\{4,6,8,9,10,12\}$ and orders $1,2,3$.
The jet logarithms are independent symbols in these finite checks; the
mathematical theorem establishes validity after arbitrary specialization.

The level-twelve polynomial check verifies both identities in
\eqref{eq:certificateq12}. The $S_4$ check constructs all 96 rows from the
two binomial rules, proves the witness equation, and computes the two ranks.
The exact run also records the first nonzero principal trace jets for six
levels. A separate finite Taylor verifier checks the next two principal
coefficients at those levels and the twisted level-260 leading coefficient.
JSON outputs record dimensions, coefficients and software versions.

The ten numerical checks use 45 decimal working digits. They include the
three level-30 traces, two twisted level-20 traces, two complex-order
conductor identities, two Stieltjes descent identities, and the retained
$S_4$ candidate. They are neither interval enclosures nor proofs of equality.
The explicit level-260 identity is proved analytically; this package does
not claim to have performed a separate 96-term high-precision evaluation
of that example.

\begin{center}
\begin{tabular}{p{0.40\textwidth}p{0.48\textwidth}}
\toprule
Artifact & Purpose\\\midrule
\code{src/distribution.py} & Explicit matrices, symbolic jet rows, level-twelve normal form.\\
\code{src/verify\_exact.py} & Exact rank and polynomial checks, with JSON output.\\
\code{src/verify\_numeric.py} & Pole-subtracted floating-point smoke tests.\\
\code{src/verify\_trace\_coefficients.py} & Finite Taylor checks for principal and twisted traces.\\
\code{src/verify\_s4\_obstruction.py} & Independent finite law-span obstruction.\\
\code{integration/} & Replacement theorem and editorial notes, not an automatic repository write.\\
\code{MANIFEST.sha256} & Integrity hashes for the distributed package.\\\bottomrule
\end{tabular}
\end{center}

No remote repository file was modified. The integration fragments require
normal editorial review, particularly comparison with any historical
matrices recovered later. The Python generator uses exact rational or
symbolic domain arithmetic for rank; there is no numerical rank tolerance.
It is a reference implementation, not a bit-complexity claim or a
production-scale character-table engine.

\section{Further research questions}

\subsection{An integral version of the polynomial normal form}

The proof uses character idempotents and division by group orders. Over a
splitting field this is harmless; over $\Z[t_p]$ it hides torsion and
integral-lattice questions. Determine an integral presentation adapted to
prime powers, and characterize the Smith data after integer-weight
specialization. Classical ordinary-distribution results suggest useful
structures, but the present field-level proof does not establish integral
freeness for this deformation.

\subsection{The conductor transition determinant}

The matrix from minimal-conductor generators to level-$q$ character sums is
diagonal with entries $A_{q,f,\ch}$. Its determinant and its zero divisor
are therefore explicit. A systematic treatment should relate the vanishing
multiplicities to arithmetic splitting of added primes, and should distinguish
coordinate singularities from intrinsic period relations. This is a good
place to develop exact conditioning estimates for character-based numerical
algorithms.

\subsection{Certified order jets at roots of unity}

Implement Proposition~\ref{prop:eval} with ball arithmetic for Stieltjes
constants and exact cyclotomic weights. Derive a precision budget that
accounts for cancellation in a trace with a prescribed high-order zero.
The finite Fourier formula supplies an exact target; the missing part is
a sharp certified numerical analysis and a tested interval implementation.

\subsection{A proof of the \texorpdfstring{$S_4$}{S4} identity outside the tested row span}

The witness in Proposition~\ref{prop:s4obstruction} excludes one route, not
the identity. A focused next step is to add weight-five higher-depth
relations and eliminate the triple coordinates, or to seek a functional
transformation of the integrals \eqref{eq:S4int}. A successful proof should
provide an exact law certificate whose new row does not annihilate the
witness. The candidate must remain explicitly conjectural until then.

\subsection{Higher-depth conductor descent}

At depth one, removing an added prime is controlled by an Euler factor.
For colored double and triple sums, prime divisibility interacts with the
strict inequalities between indices. Determine a finite operator calculus
that records these interactions and produces a conductor-adapted filtration.
A straightforward product of the depth-one factors should not be assumed
to work; the collision terms are the mathematical content of the problem.

\subsection{Formalization of the rank theorem}

A proof-assistant development can be divided into finite character
orthogonality, the two prime-edge identities, a commuting divisor-lattice
normal form, and base change to truncated germs. The rank theorem itself
requires no transcendental arithmetic. The analytic realization can be
formalized later, using the established Hurwitz multiplication and
parameter-derivative laws.

\subsection{Numerical independence remains separate}

The conductor quotient supplies exact formal coordinates, not a lower
bound for the dimension of their evaluated periods. A meaningful
independence conjecture must state the coefficient field, which lower
periods are quotiented out, and which order derivatives are included.
Special logarithmic identities and functional equations must first be
removed. Failed integer-relation searches do not fill this gap.

\section{Conclusion}

The unresolved distribution-rank pattern can be proved uniformly once the
presentation is stated and the denominator lattice is organized by minimal
character conductors. The resulting module is polynomially free, so neither
specialization nor passage to arbitrary jets creates a rank jump. This gives
an explicit theorem for the manuscript's complete distribution system and
an all-order descent formula for its Stieltjes tower.

The same structure, reflected by $s\mapsto1-s$ and character conjugation,
produces exact cyclotomic polylogarithm trace identities. At order one,
prime support controls the first nonzero derivative; the residue of the
Riemann zeta function lowers the principal vanishing order by exactly one.
These identities also expose why order-one numerical evaluation must remove
poles symbolically.

Finally, the $S_4$ certificate separates a plausible identity from an
insufficient collection of proof laws. Together, the positive normal-form
theorems and this finite obstruction give concrete next steps without
confusing formal rank, numerical evidence and arithmetic independence.

\appendix
\section{Exact normal-form algorithm in conductor coordinates}

A symbolic implementation need not perform Gaussian elimination on all
$q$ residue coordinates. The following procedure describes the mathematics
of Theorem~\ref{thm:free}.

For each divisor $d\mid q$, compute the character transform on $\U_d$.
Attach to each transformed coordinate its primitive conductor $f$.
Replace that coordinate by $A_{d,f,\ch}(\mathbf t)Z_{f,\ch}$ using
\eqref{eq:Apoly}. Collect identical primitive-character generators.
For order jets, replace $t_p$ by its truncated exponential and multiply
polynomials modulo $\eps^{N+1}$. Finally invert the finite character
transforms only for the original residue coordinates actually requested.

Every operation is finite and every transition multiplier is polynomial.
The coefficient field may be chosen as a cyclotomic field containing the
character values. Complex characters must be paired or represented exactly;
rounding them to floating-point roots of unity defeats the exactness of the
method. The reference package implements the full residue matrices and the
entire level-twelve normal form, while leaving a general optimized
character-table implementation as a subsequent engineering task.

\section{Source snapshot and claim ledger}

The repository branch was resolved during this audit to commit
\begin{center}\code{85c89e5d0bf785c058bd40fba286588e32a11382}.\end{center}
The inspected Chapter~8 blob is
\begin{center}\code{0a5b5ec8b7cb2213302c5a0b857a43e0f99b6236}.\end{center}
The audit targets are the rank observation in Chapter~8, the surviving
research programme in Chapter~10, and the $S_4$ display in Chapter~4.
The source's earlier provenance and its independently prepared research
packages are not overwritten by this note.

\begin{longtable}{p{0.49\textwidth}p{0.39\textwidth}}
\toprule
Claim & Status in this package\\\midrule
\endfirsthead
\toprule Claim & Status in this package\\\midrule\endhead
Multiparameter conductor freeness & Proved, Theorem~\ref{thm:free}.\\
Uniform endpoint-fixed rank for the complete presentation & Proved, Corollaries~\ref{cor:rank} and~\ref{cor:stielrank}.\\
Equality with an unavailable historical script's exact row list & Not asserted.\\
All-order jet dimensions & Proved, Theorem~\ref{thm:jets}.\\
All-order Stieltjes conductor identities & Proved, Theorem~\ref{thm:stieltjes}; $k=0$ correction explicit.\\
Principal and twisted trace identities & Proved, Sections~7--9.\\
Global novelty of universal-distribution rank phenomena & Not asserted; classical literature credited.\\
$S_4$ displayed closed form & Numerically supported, not proved here.\\
Obstruction in the specified $96\times24$ presentation & Exactly certified; not a disproof of $S_4$.\\
Independence of evaluated periods & Not asserted.\\
Formal verification in Lean or another proof assistant & Not performed.\\
Rigorous numerical interval certification & Not performed.\\
Remote repository integration & Not performed; fragments supplied.\\\bottomrule
\end{longtable}

\begin{thebibliography}{99}
\bibitem{repo}
V.~Reshetnikov, \emph{Polylogarithms and their Arithmetic Bridges},
ProveIt repository, consolidated manuscript, October 2026.
\url{https://github.com/VladimirReshetnikov/ProveIt/tree/85c89e5d0bf785c058bd40fba286588e32a11382/Analysis/Polylogarithms/docs/manuscript}.

\bibitem{kubert}
D.~S.~Kubert, \emph{The universal ordinary distribution},
Bulletin de la Soci\'et\'e Math\'ematique de France \textbf{107} (1979),
179--202. DOI: \href{https://doi.org/10.24033/bsmf.1891}{10.24033/bsmf.1891}.

\bibitem{ouyang}
Y.~Ouyang, \emph{On the universal norm distribution},
arXiv:math/0210297 (2002).
\url{https://arxiv.org/abs/math/0210297}.

\bibitem{coffey}
M.~W.~Coffey, \emph{Series representations for the Stieltjes constants},
Rocky Mountain Journal of Mathematics \textbf{44} (2014), 443--477;
arXiv:0905.1111. DOI: 10.1216/RMJ-2014-44-2-443.

\bibitem{dlmfhurwitz}
NIST Digital Library of Mathematical Functions,
Section~25.11, \emph{Hurwitz Zeta Function}.
\url{https://dlmf.nist.gov/25.11}. Accessed 9 October 2026.

\bibitem{dlmfpolylog}
NIST Digital Library of Mathematical Functions,
Section~25.12, \emph{Polylogarithms}.
\url{https://dlmf.nist.gov/25.12}. Accessed 9 October 2026.

\bibitem{dlmfdirichlet}
NIST Digital Library of Mathematical Functions,
Section~25.15, \emph{Dirichlet L-functions}.
\url{https://dlmf.nist.gov/25.15}. Accessed 9 October 2026.
\end{thebibliography}
\end{document}
